%! Author = zhiweizhang
%! Date = 10.06.25

\chapter{Conclusion and Future Work}\label{chapter:conclusion-and-future-work}

\section{Conclusion}

This thesis presented a comprehensive study on classifying Nazi symbols in images, exploring a wide range of
architectures—including \ac{CNN}s (ResNet and EfficientNet), transformer-based
models (\ac{ViT}), classification-specific versions of \ac{YOLO}, and multimodal approaches leveraging OpenCLIP
embeddings. In addition to end-to-end deep learning pipelines, hybrid approaches that combined OpenCLIP features
with classical \ac{ML} classifiers were also evaluated.

Two classification settings were explored: binary classification (distinguishing Nazi from non-Nazi symbols) and
multi-class classification across 13 Nazi-related symbol categories.

In the binary task, hybrid models—particularly OpenCLIP paired with classical classifiers such as \ac{SVC},
\ac{LR}, KNeighborsClassifier, and SGDClassifier—consistently delivered high performance, achieving F1-scores
above 0.89 and, in some cases, perfect \ac{ROC} \ac{AUC} scores (e.g., \ac{LR} with \ac{AUC} = 1.00). Transformer-based
\ac{ViT} and classification-oriented \ac{YOLO} also demonstrated strong results, balancing accuracy, precision,
and recall. These findings affirm the value of both pre-trained visual-semantic embeddings and end-to-end deep
feature extractors for sensitive symbol classification.

Nonetheless, some hybrid configurations—such as OpenCLIP combined with Decision Tree or Bagging Classifiers—underperformed,
especially in \ac{PR} metrics, signaling limited robustness under class imbalance. Similarly, EfficientNet
achieved high \ac{AUC} (0.993) but lower \ac{AP} (0.780), making it less suitable for applications
requiring fine-grained discrimination.

In the multi-class setting, \ac{YOLO} emerged as the top performer, achieving the highest accuracy, precision, and
F1-score. \ac{ViT} also performed strongly, confirming the effectiveness of global attention mechanisms for nuanced
symbol differentiation. While OpenCLIP performed poorly in zero-shot mode, its performance improved significantly
when combined with supervised classifiers—though still behind end-to-end models. These results highlight the
importance of downstream supervision when applying foundation models to fine-grained classification tasks.

In summary, this work demonstrates that both deep learning and hybrid approaches can be effective for the
classification of Nazi symbols. The choice of model should be guided by task constraints, desired performance
metrics, and deployment requirements. Success ultimately depends on a balance between representational power,
classifier robustness, and sensitivity to class imbalance and symbol diversity—critical considerations for building
practical and ethically responsible systems.


\section{Future Work}

Building on the encouraging results of this study, several key directions for future research emerge:

\paragraph{Dataset Expansion and Balancing:}

Increasing the diversity and quantity of labeled
images—particularly for rare or visually ambiguous Nazi symbols—could significantly improve performance in the
multi-class setting. Incorporating hard negatives (e.g., geometric symbols or cultural motifs that resemble Nazi
iconography) would also enhance binary classifier robustness against false positives.

\paragraph{Fine-Tuning Vision Backbones and Multimodal Retraining:}
While OpenCLIP embeddings paired with
classical classifiers showed strong performance, they were limited in zero-shot generalization. Task-specific
fine-tuning of the OpenCLIP vision encoder—or end-to-end multimodal retraining with labeled Nazi symbol data—may
lead to better feature separability and improved recall across all classes.

\paragraph{Precision-Focused Learning Objectives:}
As shown in the results, high \ac{AUC} did not always
translate into high real-world precision. Future work could experiment with cost-sensitive loss functions,
confidence calibration, or threshold tuning to explicitly optimize for precision-recall trade-offs—especially
important in applications with legal or reputational risks.

\paragraph{Explainability and Human Oversight:}
Integration of explainable AI (XAI) tools such as
\ac{Grad-CAM}~\parencite{Selvaraju_2019} or \ac{SHAP}~\parencite{SHAP2025} would improve transparency and
accountability. Providing visual or textual explanations for predictions could support expert review and
increase user trust—especially in archival or forensic applications.

\paragraph{Contextual and Multimodal Classification:}
Some misclassifications were driven by lack of context.
Incorporating auxiliary metadata (e.g., captions, filenames, web text) and exploring multimodal fusion models
could improve classification of ambiguous or reused symbols, and help distinguish educational versus
propagandistic usage.

\paragraph{Real-Time Deployment and System Scalability:}
Future engineering work should extend the current
containerized \ac{API} toward production-ready deployment. Key goals include low-latency inference,
cross-platform scalability, and compliance with privacy regulations (e.g., GDPR), making the system usable in
content moderation or institutional archiving workflows.

\paragraph{Human-in-the-Loop Feedback and Adaptation:}
Building mechanisms for expert feedback and iterative
improvement would support continuous learning. This is especially valuable in evolving symbol ecosystems where
new variants or context-specific uses emerge over time.


In summary, future efforts should aim to enhance the system’s precision, transparency, and adaptability to sensitive real-world environments. By combining improved data practices, targeted fine-tuning, explainability, and user input, the classifier can evolve into a more trustworthy and context-aware tool for ethical Nazi symbol recognition and moderation.

